%====================================
% KAPITEL 3: METHODIK DER THESIS
%====================================
\section{Methodik der Thesis}

\subsection{Forschungsmethodisches Vorgehen}

Zur Beantwortung der Forschungsfragen wird eine systematische Literaturrecherche durchgef\"uhrt. Dieses Vorgehen eignet sich besonders, um den bestehenden Wissensstand zu Change-Management-Faktoren bei ERP-Implementierungen im Mittelstand zu erfassen, zu strukturieren und kritisch zu synthetisieren. Die systematische Literaturrecherche folgt den etablierten Rahmenwerken von Webster und Watson\footcite[Vgl.][S.~xiii--xxiii]{webster2002}, vom Brocke et al.\footcite[Vgl.][]{vombrocke2015} sowie Okoli\footcite[Vgl.][]{okoli2015}, die einen nachvollziehbaren und reproduzierbaren Forschungsprozess sicherstellen.

Die Literatursuche erfolgt in f\"unf wissenschaftlichen Datenbanken, die f\"ur die Wirtschaftsinformatik und angrenzende Disziplinen relevant sind: SpringerLink, ScienceDirect, IEEE Xplore, EBSCOhost sowie die AIS Electronic Library (AISeL). Der Suchalgorithmus kombiniert Boolesche Operatoren, um die zentralen Konstrukte der Arbeit pr\"azise abzudecken. Die Suchstring-Formulierung lautet:

\begin{quote}
\verb|("ERP implementation" OR "Enterprise Resource Planning")|\\
\verb|AND ("change management")|\\
\verb|AND ("critical success factors" OR "success factors")|\\
\verb|AND ("SME" OR "Mittelstand")|
\end{quote}

F\"ur die deutschsprachige Suche wird eine \"aquivalente Variante verwendet:

\begin{quote}
\verb|("ERP-Einf"uhrung" OR "ERP-Implementierung")|\\
\verb|AND ("Change Management" OR "Ver"anderungsmanagement")|\\
\verb|AND ("Erfolgsfaktoren" OR "Mittelstand" OR "KMU")|
\end{quote}

Die Qualit\"atssicherung der identifizierten Literatur erfolgt \"uber mehrere Stufen. Ber\"ucksichtigt werden ausschlie\ss{}lich peer-reviewte Beitr\"age aus wissenschaftlichen Zeitschriften und Konferenzb\"anden sowie etablierte Fachb\"ucher. Die Zeitschriftenqualit\"at wird anhand des H-Index und der Einstufung im VHB-JOURQUAL-Ranking bewertet. Das angestrebte Verh\"altnis zwischen Journalartikeln und Fachbuchbeitr\"agen liegt bei etwa 60 zu 40, um sowohl aktuelle empirische Befunde als auch theoretische Fundamentierung zu gew\"ahrleisten.

Als theoretischer Rahmen dienen zwei etablierte Modelle. Das IS-Erfolgsmodell nach DeLone und McLean\footcite[Vgl.][S.~9--30]{delone2003} operationalisiert den Erfolg von Informationssystemen \"uber die Dimensionen Systemqualit\"at, Informationsqualit\"at, Servicequalit\"at, Nutzerzufriedenheit und Nettonutzen. Das ADKAR-Modell\footcite[Vgl.][S.~124--132]{galli2018} strukturiert Change-Management-Ma\ss{}nahmen entlang der Stufen Awareness, Desire, Knowledge, Ability und Reinforcement. Die Verkn\"upfung beider Modelle erfolgt \"uber die literaturbasiert hergeleitete Annahme, dass die ADKAR-Phasen auf die D\&-M-Erfolgsdimensionen wirken: Awareness und Desire beeinflussen die Nutzerakzeptanz und damit die User Satisfaction; Knowledge und Ability steigern die Systemnutzungseffizienz und den Individual Impact; Reinforcement sichert die langfristige Systemnutzung und verst\"arkt die Net Benefits. Diese Kombination erlaubt es, sowohl die technisch-organisationale als auch die personelle Perspektive der ERP-Implementierung abzubilden.

Die Auswertung der identifizierten Faktoren erfolgt mittels einer inhaltlich strukturierenden Inhaltsanalyse nach Mayring und Fenzl.\footcite[Vgl.][S.~633--648]{mayring2019} Dabei werden die in der Literatur genannten Change-Management-Faktoren zun\"achst extrahiert, anschlie\ss{}end induktiv und deduktiv zu Kategorien verdichtet und schlie\ss{}lich hinsichtlich ihrer Relevanz f\"ur den Mittelstand bewertet.

\subsection{Untersuchungsgegenstand und Datenbasis}

Gegenstand der Untersuchung ist die aktuelle Forschungsliteratur zu Change-Management-Faktoren bei ERP-Implementierungen, mit einem Schwerpunkt auf Publikationen ab dem Jahr 2015. Diese zeitliche Eingrenzung stellt sicher, dass der technologische und organisatorische Wandel der vergangenen Jahre -- insbesondere die zunehmende Cloud-Nutzung und agile Implementierungsmethoden -- in den analysierten Studien abgebildet wird.

Die angestrebte Datenbasis umfasst voraussichtlich etwa 20 bis 30 peer-reviewte Quellen aus wissenschaftlichen Zeitschriften und Konferenzb\"anden. Davon werden mindestens drei als Kernstudien identifiziert, die sich spezifisch mit Change-Management-Faktoren im Kontext kleiner und mittlerer Unternehmen befassen. Die geografische Eingrenzung auf den deutsch- und englischsprachigen Raum gew\"ahrleistet sowohl die Vergleichbarkeit der wirtschaftlichen Rahmenbedingungen als auch die Breite des internationalen Forschungsstands. Um die Transparenz des Auswahlprozesses zu erh\"ohen, wird ein PRISMA-Flowchart erstellt, das die Anzahl der identifizierten, gescreenten und final eingeschlossenen Publikationen dokumentiert.\footcite[Vgl.][Art.-Nr.~n71]{page2021}

Aus der systematischen Synthese der Literatur wird erwartet, dass sich signifikante Change-Management-Hebel identifizieren lassen. Dazu z\"ahlen insbesondere eine transparente und kontinuierliche Kommunikation w\"ahrend des gesamten Implementierungsprozesses, die aktive Einbindung und Schulung der betroffenen Mitarbeiter sowie ein klares Commitment und Vorbildverhalten der F\"uhrungsebene. Diese Faktoren werden im anschlie\ss{}enden Kapitel detailliert analysiert und hinsichtlich ihrer spezifischen Auspr\"agung im Mittelstand diskutiert.
